% ---
% id: reductions/wpfss/def:targets
% kind: def
% title: "Certified weights and the flag-state and weight-preserving targets"
% summary: >-
%   Certified weights are per-outcome lower bounds on a POVM's blocks above a
%   cutoff. Both targets keep the blocks below the cutoff; the flag-state target
%   puts each outcome on its own flag, the weight-preserving target puts its
%   certified weight on every flag and one minus the total weight on its own.
% prominence: standard
% status: accepted
% uses: [def:weight-route]
% mentions: [prop:chain, rem:published]
% assumes: []
% refs: []
% provenance: newly-derived
% lean:
% ---

Above a photon-number cutoff the receiver is replaced by one classical flag
per outcome, and the only device knowledge kept there is a lower bound on each
outcome's block. The weight-preserving target keeps these bounds in its flags,
so that the corresponding inequality holds on every squashed state
(Proposition~\ref{prop:chain}\,(iii)); both targets are known from the
literature (Remark~\ref{rem:published}).

\begin{definition}
\label{def:targets}
Throughout this definition and the claims of this submodule that use it,
$\Hs{\Bsys}=\Hs{\Bsys_{\leqN}}\dsum\Hs{\Bsys_{\gtN}}$
is an orthogonal two-block decomposition with projections $\Pleq$ and $\Pgt$
--- a photon-number cutoff in the optical instance, where
$\Hs{\Bsys_{\leqN}}=\bigoplus_{m\le N}\Hs{\Chan_m}$ is finite-dimensional ---
and every $M_v$ is block diagonal for it,
$M_v=\blk{M}{\leqN}_v\dsum\blk{M}{\gtN}_v$. $\Hs{\Flag}=\mathbb{C}^{|\Vset|}$
carries the orthonormal basis $\{\ket{v}\}_{v\in\Vset}$, and
$\Hs{\Qsys}=\Hs{\Bsys_{\leqN}}\dsum\Hs{\Flag}$ with block projections
$\Pleq$ (the same symbol on either space) and $\Pflag$.

A family $(p_v)_{v\in\Vset}$ of
\term[Non-negative numbers, one per outcome, each an operator lower bound on that outcome's POVM element above the cutoff.]{certified weights}
for $M$ at the cutoff consists of non-negative numbers $p_v$ --- on the
emission path, integer multiples of $2^{-\kappa}$, the $p_{(y,t)}$ of
\eqref{eq:weightroute} --- with
\begin{equation}\label{eq:certweights}
  \blk{M}{\gtN}_v\;\succeq\;p_v\,\idon{\Bsys_{\gtN}}
  \qquad\text{for every }v\in\Vset ,
\end{equation}
and we write $p\coloneqq\sum_vp_v$ for their
\term[The sum of the certified weights over all outcomes.]{total weight}.
Define
\begin{align}
  \Ffss_v&\coloneqq\blk{M}{\leqN}_v\dsum\dyad{v}_{\Flag},
  \label{eq:FSS}\\
  \Fwp_v&\coloneqq\blk{M}{\leqN}_v\dsum
     \big(p_v\,\idon{\Flag}+(1-p)\,\dyad{v}_{\Flag}\big).
  \label{eq:WPFSS}
\end{align}
$\Ffss$ is the
\term[The target that keeps the POVM below the cutoff and replaces everything above it by one flag per outcome: the target of the flag-state squasher.]{flag-state target}
and $\Fwp$ the
\term[The target that keeps the POVM below the cutoff and gives each outcome a flag block carrying its certified weight on every flag and, in addition, one minus the total weight on its own flag: the target of the multi-bound weight-preserving flag-state squasher.]{weight-preserving target},
the multi-bound weight-preserving flag-state-squasher target, with one
certified weight per outcome. The
\term[The weight-preserving target in which only one certified weight is non-zero.]{single-bound case}
is the one in which only one $p_v$ is non-zero, and $\Ffss$ is the case
$p_v\equiv0$.
\end{definition}
